\chapter{Continuum Mechanics Caveats}
\label{app:continuum-caveats}

The continuum chapters use compact notation because the main text already
derives the equations in detail.  This appendix gathers the assumptions that
are easiest to blur: what is kinematics, what is balance law, what is
constitutive law, which stress measure is being used, and where idealizations
enter.  The goal is to prevent correct formulas from being read with the wrong
meaning.

\section{Three Layers: Kinematics, Balance, Constitutive Law}
\label{sec:continuum-three-layers}

A continuum theory is built in three logically distinct layers.

\begin{enumerate}
\item \emph{Kinematics} specifies the variables.  For a material body this may
be a deformation map \(\varphi(X,t)\), deformation gradient
\(F=\pd\varphi/\pd X\), velocity \(v(x,t)\), or director field.
\item \emph{Balance laws} express conservation of mass, momentum, and angular
momentum.  They hold for broad classes of materials before a particular
stress-strain relation has been chosen.
\item \emph{Constitutive laws} specify the material response, such as
\(\sigma=-pI+2\mu D\) for a Newtonian incompressible fluid or
\(P=\pd W/\pd F\) for a hyperelastic solid.
\end{enumerate}

Most mistakes in continuum mechanics come from mixing these layers.  For
example, Cauchy's momentum equation
\[
        \rho\frac{D v}{D t}=\divergence\sigma+\rho b
\]
is a balance law.  The Navier-Stokes equation appears only after inserting a
Newtonian constitutive law and, for the incompressible form, the constraint
\(\divergence v=0\).

\section{Reference And Spatial Descriptions}
\label{sec:reference-spatial-descriptions}

Let \(\mathcal{B}\) be a reference body with material point \(X\), and let
\(\varphi_t:\mathcal{B}\to\R^3\) be the placement at time \(t\).  The spatial
point is
\[
        x=\varphi(X,t).
\]
The deformation gradient is
\[
        F^i{}_A=\frac{\pd \varphi^i}{\pd X^A}.
\]
Here Latin indices \(i,j,\dots\) refer to spatial components, while capital
indices \(A,B,\dots\) refer to material components.  The Jacobian
\[
        J=\det F
\]
measures local volume change.  If \(\rho_0(X)\) is the reference mass density
and \(\rho(x,t)\) is the spatial density, conservation of mass gives
\[
        \rho(\varphi(X,t),t)J(X,t)=\rho_0(X).
\]

\begin{figure}[htbp]
\centering
\begin{tikzpicture}[scale=1.0, >=Stealth]
\draw[thick] (-2.8,-0.9) .. controls (-2.2,-1.4) and (-1.1,-1.0) .. (-0.7,-0.2)
    .. controls (-0.4,0.5) and (-1.0,1.1) .. (-2.0,1.0)
    .. controls (-3.0,0.8) and (-3.3,-0.3) .. (-2.8,-0.9);
\draw[thick] (2.0,-1.0) .. controls (3.0,-1.3) and (4.0,-0.7) .. (4.2,0.2)
    .. controls (4.4,1.0) and (3.4,1.5) .. (2.3,1.1)
    .. controls (1.4,0.8) and (1.1,-0.4) .. (2.0,-1.0);
\fill (-2.0,0.2) circle (1.5pt) node[above left] {\(X\)};
\fill (2.9,0.35) circle (1.5pt) node[above right] {\(x=\varphi(X,t)\)};
\draw[->, thick] (-0.4,0.25) -- (1.55,0.25) node[midway, above] {\(\varphi_t\)};
\draw[->] (-2.0,0.2) -- (-1.4,0.55) node[midway, above] {\(\delta X\)};
\draw[->] (2.9,0.35) -- (3.72,0.78) node[midway, above] {\(F\delta X\)};
\node at (-2.0,-1.55) {reference body \(\mathcal{B}\)};
\node at (2.95,-1.55) {current body \(\varphi_t(\mathcal{B})\)};
\end{tikzpicture}
\caption{Material and spatial descriptions are connected by the placement
\(\varphi_t\).  The deformation gradient \(F\) maps material tangent vectors
to spatial tangent vectors.}
\label{fig:continuum-reference-spatial-map}
\end{figure}

\section{Stress Measures}
\label{sec:stress-measures-table}

Different stress tensors answer different questions.

\begin{center}
\begin{tabular}{p{0.18\textwidth} p{0.28\textwidth} p{0.43\textwidth}}
\hline
symbol & name & meaning\\
\hline
\(\sigma\) & Cauchy stress & spatial traction per current area:
\(t=\sigma n\)\\
\(P\) & first Piola-Kirchhoff stress & spatial force per reference area;
appears in \(\operatorname{Div}P+\rho_0 b=\rho_0\ddot\varphi\)\\
\(S\) & second Piola-Kirchhoff stress & material stress work-conjugate to
\(E=(C-I)/2\), where \(C=F^TF\)\\
\(\tau=J\sigma\) & Kirchhoff stress & volume-weighted Cauchy stress, useful in
finite strain\\
\hline
\end{tabular}
\end{center}

The standard Piola transform is
\[
        P=J\sigma F^{-T}.
\]
If the material is hyperelastic with stored energy density \(W(F)\) per
reference volume, then
\[
        P=\frac{\pd W}{\pd F}.
\]
Objectivity requires
\[
        W(RF)=W(F)\qquad \text{for all }R\in\SO(3),
\]
so that the energy does not change under a superposed rigid rotation.

\section{Incompressibility And Pressure}
\label{sec:continuum-incompressibility-pressure}

Incompressibility means different things in material and spatial language.  In
finite elasticity it is
\[
        J=\det F=1.
\]
In Eulerian fluid mechanics it is
\[
        \divergence v=0.
\]
The second condition follows from differentiating the first along a material
trajectory:
\[
        \dot J=J\,\divergence v.
\]
In both settings, pressure is a Lagrange multiplier for a constraint.  For an
incompressible hyperelastic solid one writes
\[
        P=\frac{\pd W}{\pd F}-pJF^{-T}.
\]
For an incompressible Newtonian fluid one writes
\[
        \sigma=-pI+2\mu D,\qquad
        D=\frac{1}{2}\bigl(\nabla v+\nabla v^T\bigr).
\]
The pressure is determined together with the velocity by the constraint and
the boundary conditions; it is not supplied by an equation of state in the
incompressible model.

\section{Elastic Solids Versus Fluids}
\label{sec:elastic-versus-fluid-caveats}

The notes place elastic bodies before deforming bodies and fluids because
elasticity makes the material reference configuration explicit.  Fluids can
also be described materially, but their standard equations are Eulerian and
their constitutive law does not remember a fixed rest shape.

\begin{center}
\begin{tabular}{p{0.27\textwidth} p{0.29\textwidth} p{0.29\textwidth}}
\hline
feature & elastic solid & Newtonian fluid\\
\hline
primary memory & deformation from reference body & instantaneous rate of
strain\\
typical constitutive variable & \(F\), \(C=F^TF\), or strain \(E\) &
\(D=(\nabla v+\nabla v^T)/2\)\\
stress response & elastic stress from stored energy & viscous stress plus
pressure\\
equilibrium equation & boundary-value problem for displacement & steady or
unsteady velocity-pressure problem\\
constraint pressure & enforces \(J=1\) when incompressible & enforces
\(\divergence v=0\) when incompressible\\
\hline
\end{tabular}
\end{center}

The distinction is physical, not merely notational.  A solid can sustain shear
stress at rest.  A Newtonian fluid at rest has stress \(-pI\); shear stress is
proportional to rate of deformation.

\section{Boundary Conditions Are Part Of The Model}
\label{sec:continuum-boundary-conditions}

The same differential equation can describe different physical problems under
different boundary conditions.  For elasticity, common boundary conditions are
prescribed displacement
\[
        u=u_0\quad \text{on }\Gamma_D
\]
and prescribed traction
\[
        P N=T_0\quad \text{on }\Gamma_N .
\]
For a viscous fluid at a fixed wall, the usual no-slip condition is
\[
        v=v_{\rm wall}.
\]
For inviscid flow, only impermeability
\[
        v\cdot n=v_{\rm wall}\cdot n
\]
is normally imposed.  Applying a no-slip condition to the Euler equation or an
inviscid slip condition to a viscous wall problem changes the physical model.

\section{Navier-Stokes Assumptions}
\label{sec:navier-stokes-assumptions}

The incompressible Navier-Stokes equation
\[
        \rho(\pd_t v+v\cdot\nabla v)
        =
        -\nabla p+\mu\Delta v+\rho b,\qquad
        \divergence v=0,
\]
uses the following assumptions:
\begin{enumerate}
\item continuum approximation: fields are smooth enough on the length scales
of interest;
\item constant density \(\rho\) and viscosity \(\mu\);
\item Newtonian, isotropic, frame-indifferent viscous stress;
\item incompressibility;
\item boundary and initial data compatible with \(\divergence v=0\).
\end{enumerate}
The equation is local and deterministic.  Turbulence does not invalidate it as
a model, but turbulent prediction may require statistical descriptions,
resolution studies, or closure assumptions if one averages the equations.

\section{Energy Identities And Regularity}
\label{sec:continuum-energy-regularity}

For smooth incompressible flow on a domain with boundary conditions that remove
boundary power, multiplying Navier-Stokes by \(v\) and integrating gives
\[
        \frac{\dd}{\dd t}\int \frac{1}{2}\rho |v|^2\,\dd x
        =
        -\int 2\mu D:D\,\dd x
        +
        \int \rho b\cdot v\,\dd x .
\]
The pressure term drops out because
\[
        \int v\cdot\nabla p\,\dd x
        =
        \int \divergence(pv)\,\dd x-\int p\,\divergence v\,\dd x.
\]
The second term is zero by incompressibility, and the first is a boundary flux.

For weak solutions, the same calculation may become an inequality rather than
an equality.  In three dimensions, global smoothness for arbitrary smooth data
is a famous open problem.  The lecture-note derivations are therefore
classical derivations: they are correct under the stated smoothness
assumptions, while the functional-analytic theory requires additional care.
